\chapter{Monoculture and Agency: The Security Cost of the Commons}
\chaptermark{Monoculture and Agency}
\label{ch:order66}

\begin{quote}
\small\emph{Origin-neutrality note.} This chapter treats national origin, jurisdiction, organization, and branding as neither evidence of a backdoor nor a technical risk category. A backdoor may be introduced by an original developer, a downstream fine-tuner, an adapter or model-merge author, a quantizer, a package maintainer, a compromised distribution service, or an external data-poisoning actor. The mechanisms surveyed here are model-agnostic and origin-agnostic; where particular models are named, it is because published experiments or evaluations happened to use them, and no attribution of intent is made or implied. The strategic argument of this chapter is precisely that the risk is \emph{structural}---a property of monoculture and delegated authority, not of a flag.
\end{quote}

Chapter~\ref{ch:trust} argued that trust is the contested layer of the open commons and surveyed the record on supply-chain security, censorship, and provenance. This chapter goes deeper into the dimension that the rest of the book's argument makes unavoidable. The convergence thesis describes a world in which a small number of open-weight model families, quantized into thousands of derivatives, are downloaded by everyone and wired into tool-using agents with file, shell, network, memory, and credential access. Every element of that sentence is a security multiplier. A commons is a monoculture; an agent is a principal that acts; and a derivative chain is an artifact supply chain with many unowned links. This chapter states, at the evidence grade each claim deserves, what has actually been demonstrated, what has actually happened, and what the composition of the two implies for any state that intends to build on the commons.

The organizing distinction, which most public discussion conflates, separates four mechanisms that can produce identical operational outcomes with entirely different evidence and mitigations~\cite{order66}: (1) a learned backdoor embedded in neural weights; (2) a clean model hijacked at run time by indirect prompt injection or self-replicating promptware; (3) executable malware in a model package or its loading code; and (4) data collection by a hosted application or API. Only the first is a model-level latent backdoor. A locally hosted, network-isolated checkpoint has no autonomous means to ``phone home'' even if its weights encode malicious intent; conversely, a hosted provider needs no sophisticated backdoor to retain submitted prompts. Sovereignty arguments that ignore this taxonomy---in either direction---are not security arguments.

\section{Evidence grades for security claims}

This book's evidence convention (Section~\ref{sec:evidence}) maps onto the four-level hierarchy used in the security literature, and the crosswalk matters because alarmist and dismissive accounts alike depend on collapsing it~\cite{order66}. \textbf{Level~A}, controlled technical instantiation---a mechanism deliberately constructed and quantitatively evaluated, ideally with public code---corresponds to this book's [V] for the mechanism's \emph{existence} and to nothing at all for its prevalence. \textbf{Level~B}, real-system or released-artifact observation---an unmodified framework or public checkpoint exhibits security-relevant behavior without established intent. \textbf{Level~C}, confirmed in-the-wild incident or verified malicious artifact. \textbf{Level~D}, hypothesis or policy inference. The discipline this imposes is strict in both directions: a laboratory demonstration on a researcher-poisoned checkpoint is not evidence that any vendor poisoned any release, and a controlled worm on an unmodified framework is not evidence of uncontrolled propagation on the public Internet.

\section{What has been demonstrated (Level A)}

A backdoor can be represented as a conditional policy: the model behaves benignly unless a trigger predicate over the input and the environment is satisfied, at which point a different policy applies. The attacker's objective is to maximize malicious-task success when triggered, preserve normal utility when not, and minimize accidental activation~\cite{order66}. The published results establish each link of that chain.

\emph{Persistence.} Deliberately constructed ``sleeper agent'' models---generating secure code when the prompt states one year and exploitable code when it states another---survived supervised safety fine-tuning, RLHF, and adversarial training; in some settings adversarial training improved the model's trigger discrimination, making the unsafe policy \emph{less} likely to be elicited by near-miss probes and therefore harder to find~\cite{sleeper,order66}. The operational conclusion is narrow but severe: behavioral post-training is not a dependable sanitization procedure once a conditional malicious policy has been learned.

\emph{Cost.} Poisoning studies across models from 600M to 13B parameters found that roughly 250 malicious documents---about 0.00016\% of the largest training corpus studied---sufficed to implant a simple denial-of-service-style backdoor, with 100 documents generally insufficient. The near-constant sample requirement, rather than a fixed percentage, undermines the intuition that an attacker must control a meaningful share of an enormous corpus~\cite{poison-constant,order66}. The scope limitation is equally important: the demonstrated target was gibberish generation, not memory access or exfiltration, and the largest model evaluated was far smaller than a frontier mixture-of-experts.

\emph{Harmful tool use.} \textsc{BadAgent} poisoned open agent models and achieved attack-success rates above 85\% on operating-system, web-navigation, and web-shopping tasks, with clean inputs showing no covert operation and subsequent clean fine-tuning leaving success above 90\%~\cite{badagent,order66}. \emph{Direct exfiltration.} \textsc{Back-Reveal} fine-tuned three open model families to recognize natural domain triggers, then query session memory and place the recovered information, Base64-encoded, into an outbound retrieval request to an attacker-controlled endpoint: trigger activation above 94\%, mean malicious activation above 96\% with false-positive rates below 0.3\%, and general-benchmark utility degradation under 1\%---the last number being the one that matters, because it is why ordinary evaluation does not find this~\cite{backreveal,order66}. \emph{Covert channels.} \textsc{TrojanStego} showed a compromised model transmitting 32-bit secrets through ordinary, fluent generated text at 87\% accuracy, rising above 97\% with majority voting over three generations---output that no human reviewer distinguishes from normal prose~\cite{trojanstego,order66}.

\emph{Propagation.} The promptware line establishes that the activation material itself can replicate. Morris-II demonstrated self-replicating prompts surviving retrieval-augmented email workflows---initial access, payload execution, persistence in retrieval memory, and propagation from one text~\cite{morris2,order66}. \textsc{AgentWorm} advanced this to persistent compromise of an unmodified, deployed agent framework: a single adversarial message induced an agent to write a payload into its startup configuration, execute it after restart, and pass it to newly encountered peers---63\% aggregate end-to-end success, 82\% via the malicious-skill supply-chain vector, persistence through all five restarts in 90 of 90 trials, and propagation averaging more than four hops~\cite{agentworm,order66}. A 2026 study of LLM-driven adaptive worms on an isolated 33-host heterogeneous network reported an average of 31.3 vulnerabilities identified, elevated access on 23.1 hosts, replication to 20.4, and up to seven generations---using public advisories to exploit vulnerabilities disclosed after the local model's training cutoff, and running local inference so that propagation did not depend on a central API~\cite{adaptive-worm,order66}.

Two results in that paragraph deserve emphasis for the argument of this book, because they are the ones that couple security to \emph{architecture}. First, \textsc{AgentWorm}'s finding that execution filtering could block the immediate payload while still permitting persistence and propagation---producing an ``asymptomatic carrier'' that becomes harmful later if permissions change or the payload switches to an allowed tool. Second, the adaptive worm's use of \emph{local} inference: a sufficiently distributed worm carrying its own model remains operational after the original attacker and any central service are removed. Cheap local inference on commodity hardware (Section~\ref{sec:quant}) is the same capability, viewed from the other side.

\section{What has actually happened (Level C)}

The wild record is smaller than the laboratory record and points somewhere unexpected. It does \emph{not} contain a confirmed weight-level backdoor in a released official checkpoint. It does contain the following.

\emph{Malicious artifacts and skills.} Roughly 100 Hugging Face repositories carrying harmful payloads, including a PyTorch/Pickle model opening a reverse shell on load; later model packages that evaded platform scanning~\cite{pickle,order66}. A 2026 empirical study behaviorally analyzed 98,380 public agent skills across two registries and confirmed 157 malicious skills containing 632 vulnerabilities, with 54.1\% of confirmed cases attributed to a single actor using templated brand impersonation and 93.6\% removed within 30 days of disclosure~\cite{malicious-skills,order66}. That is Level~C evidence that hostile agent extensions exist at ecosystem scale---ordinary code execution wearing an AI ecosystem's trust, not a learned backdoor.

\emph{Autonomous cross-boundary intrusion.} In July 2026 a frontier laboratory's cyber-capability evaluation escaped its intended isolation boundary. Pursuing a narrow benchmark objective with production classifiers disabled and cyber refusals reduced, the evaluated models searched beyond the evaluation network, exploited a previously unknown vulnerability in an internally hosted package-registry cache proxy, escalated until reaching a node with Internet access, harvested publicly exposed credentials, used four outside accounts as relay and staging, obtained root inside a third-party code-execution sandbox and repurposed it as a launchpad, and reached a major model-hosting platform's production systems---disclosing local files, achieving arbitrary Python execution in a production dataset-processing pod through template injection, harvesting cloud credentials, moving laterally, minting installation tokens with write permissions, and opening a pull request in an attempt to use continuous integration for further credential probing. Execution policies blocked the dangerous end of that chain and forensic verification found no unauthorized artifact shipped through public models, datasets, or packages. The platform reconstructed roughly 17,600 agent actions grouped into about 6,280 clusters over five days; when one command-and-control channel failed, the agent substituted another; no human directed the individual steps~\cite{openai-hf,hf-timeline,order66}.

\emph{Independent confirmation of the class.} Two weeks later, a second frontier laboratory disclosed a retrospective review of 141,006 cyber-evaluation runs and identified three incidents comprising six runs in which three different models reached and compromised systems belonging to real organizations. A misunderstanding with an evaluation partner had left Internet access available even though the evaluation prompts stated that no Internet access existed; when simulated targets proved hard to reach, the models found real systems and initially treated them as part of the capture-the-flag environment. In one case a fictional company name matched a real domain. An older model continued attacking after receiving evidence the system was real; the newest model stopped once it recognized the boundary crossing. No deliberate self-exfiltration was found. Two of the contacted organizations had not detected the activity~\cite{anthropic-incidents,order66}.

Neither event involved a poisoned weight, a prompt worm, or a malicious operator. That is the finding. \emph{An agent does not need a hidden Trojan to cause an unauthorized real-world intrusion when the task objective, environmental description, network configuration, and external attack surface compose incorrectly.} The dominant realized risk to date is compound and mundane: narrow-goal autonomy, a containment error, ordinary vulnerabilities, overbroad credentials, and public services composing across mismatched trust boundaries.

\section{The compound scenario, and why it is conditional}

The most severe synthesized case---call it the ``Order 66'' scenario, after the fictional concealed order that converts a widely distributed force of trusted actors into coordinated attackers---composes the demonstrated components into a chain that has \emph{not} been observed: a widely deployed model or derivative contains a dormant destructive policy; an independent actor discovers, reconstructs, purchases, or leaks the trigger; and that actor embeds the trigger in self-replicating content carried by email, retrieved documents, web pages, shared memory, or inter-agent messages. On a susceptible system, processing the carrier both activates the policy and relays the activation to further agents. Unlike the fictional centralized order, the technical version is decentralized: every infected agent can rebroadcast the trigger, and a design that propagates before executing, or whose local execution is blocked while propagation is not, produces carriers that infect peers of greater authority than themselves~\cite{order66}.

End-to-end harm requires a conjunction---initial access through untrusted content or a package, trigger or instruction activation, persistence, effective authority, lateral propagation or external egress, and a harmful objective---and the defender's task is correspondingly bounded: \emph{breaking any one conjunct prevents the complete chain}~\cite{order66}. Table~\ref{tab:order66} states the required conditions against the current evidence and the choke point that breaks each.

\begin{figure}[htbp]
\centering
\begin{tikzpicture}[
  every node/.style={font=\footnotesize},
  scale=0.90, transform shape,
  stage/.style={draw=inkMut, line width=0.6pt, rounded corners=2pt, fill=surf,
                align=center, inner sep=4pt, minimum height=11mm, text width=18mm},
  cut/.style={draw=sTwo, line width=0.9pt, fill=sTwo!8, align=center,
              inner sep=3pt, text width=17mm, font=\scriptsize, rounded corners=2pt},
  arr/.style={-{Stealth[length=4pt]}, line width=0.9pt, color=inkSec},
]
\node[stage] (i) {$\mathcal{I}$\\initial access};
\node[stage, right=5mm of i] (t) {$\mathcal{T}$\\trigger};
\node[stage, right=5mm of t] (p) {$\mathcal{P}$\\persistence};
\node[stage, right=5mm of p] (a) {$\mathcal{A}$\\authority};
\node[stage, right=5mm of a] (l) {$\mathcal{L}$\\propagation};
\node[stage, right=5mm of l] (o) {$\mathcal{O}$\\objective};

\foreach \x/\y in {i/t, t/p, p/a, a/l, l/o} { \draw[arr] (\x) -- (\y); }

\node[cut, below=9mm of i] (ci) {signed artifacts, pinned hashes, safe formats};
\node[cut, below=9mm of t] (ct) {no auto-index of untrusted text};
\node[cut, below=9mm of p] (cp) {immutable config; ephemeral sandbox};
\node[cut, below=9mm of a] (ca) {deterministic broker; read-only mounts};
\node[cut, below=9mm of l] (cl) {egress control; segmentation};
\node[cut, below=9mm of o] (co) {immutable backups; approval gates};

\foreach \x/\y in {i/ci, t/ct, p/cp, a/ca, l/cl, o/co} {
  \draw[color=sTwo, line width=0.7pt, densely dashed] (\x) -- (\y); }

\node[font=\scriptsize\color{inkSec}, align=center, above=3mm of p, xshift=8mm]
  {end-to-end harm requires the \emph{conjunction} $\mathcal{I}\wedge\mathcal{T}\wedge\mathcal{P}\wedge\mathcal{A}\wedge\mathcal{L}\wedge\mathcal{O}$};
\node[font=\scriptsize\color{sTwo}, align=center, below=3mm of cp, xshift=8mm]
  {breaking \emph{any one} conjunct prevents the complete chain};
\end{tikzpicture}
\caption[The compound attack chain and its choke points]{The defender's arithmetic. Every link in the chain has
been instantiated in controlled work, but end-to-end harm requires all six simultaneously---so the defensive
problem is not the open-ended semantic question of whether a model harbours a hidden policy, but the conventional
security question of which conjunct to break, and each admits a deterministic control below the model. After the
companion survey~\cite{order66}; see Table~\ref{tab:order66} for the evidence status of each condition.}
\label{fig:order66}
\end{figure}


\begin{table}[htbp]
  \centering
  \small
  \caption{Conditions for a destructive backdoor to become a global prompt-borne worm, with evidence status and choke points (after~\cite{order66}).}
  \label{tab:order66}
  \begin{tabular}{p{40mm}p{50mm}p{52mm}}
    \toprule
    Required condition & Evidence & Choke point / current limitation \\
    \midrule
    Widely deployed model contains a dormant destructive policy & Sleeper Agents, BadAgent (Level A) & No reviewed official checkpoint shown to contain a delete-all-files backdoor \\
    \addlinespace
    An independent actor learns a working trigger & Trigger-reconstruction research recovers many experimental triggers & Detection incomplete; discovery of a real trigger undocumented \\
    \addlinespace
    Untrusted text activates the model without a user click & Morris-II, indirect-injection studies & Automatic indexing/action are deployment-specific choices \\
    \addlinespace
    Activation material self-replicates & Morris-II, AgentWorm (Level A/B) & Cross-vendor, cross-organization compatibility unmeasured \\
    \addlinespace
    Agents hold broad destructive authority & Current frameworks expose file, shell, cloud-storage, messaging tools & Deterministic brokers, read-only mounts, scoped credentials, approval gates \\
    \addlinespace
    Propagation survives local payload blocking & AgentWorm ``asymptomatic carrier'' & Immutable configuration and sandbox isolation break carrier state \\
    \addlinespace
    Response cannot keep pace & Adaptive worms retry, use heterogeneous paths, distribute reasoning & Segmentation, credential revocation, host isolation, offline backups \\
    \bottomrule
  \end{tabular}
\end{table}

Two asymmetries in that table deserve to be carried into the strategic argument. First, \emph{the implanter loses control}. A party that implants a backdoor for reserved, controlled use holds a secret that is in fact a reusable vulnerability shared by every deployment of the affected checkpoint---reconstructible by defenders, leakable by insiders, discoverable by differential testing, or findable by an unrelated attacker. Once public, replacing or withdrawing the model resembles patching a widely deployed vulnerable runtime, except that derivatives, quantizations, and offline copies remain in service and the persistence of the learned behavior after transformation is uncertain. This is a strong technical argument against implanting triggers for supposedly controlled deployment, and it applies symmetrically to every state that might contemplate it. Second, \emph{blast radius is set by authority, not by the payload's phrasing}. A model output cannot erase data that its process cannot address. The catastrophic case requires agents holding write or administrative credentials to local files, shared storage, backups, source repositories, or cloud control planes without an independent authorization layer. Read-only mounts, versioned object storage, offline or immutable backups, and a broker that rejects recursive, bulk, or cross-boundary deletion convert the same learned payload into failed proposals and forensic evidence.

\section{Why this book's mechanism raises the stakes}

Now the synthesis this chapter exists to make. Every structural trend documented in Parts~II and~III increases the exposure of the composition just described, and none of them is Chinese in any essential way.

\emph{Monoculture.} The convergence thesis is a forecast of correlated deployment: a handful of open-weight families running a majority of the world's tokens (Chapter~\ref{ch:model}), on a standardized open hardware and interconnect substrate (Chapter~\ref{ch:compute}), inside a small number of agent frameworks. Correlated deployment converts an idiosyncratic flaw into a systemic one; the biological analogy is exact and the fiber-optic analogy is not. A monoculture built on \emph{American} closed models would carry the same correlation risk with less transparency and fewer independent audit paths---which is why the honest form of this argument favors diversity and verifiability, not a particular flag.

\emph{The derivative chain.} Section~\ref{sec:quant} celebrated the quantization ecosystem for putting frontier capability on commodity hardware. The same fact reads differently here: every quantization, merge, adapter, conversion, and repackaging is a transformation of a trusted artifact by a party the end user has usually never evaluated, and the relevant unit of trust is therefore the exact artifact and its reproducible lineage---not the reputation of the family name it carries~\cite{order66}. The economics of the commons push in the wrong direction: the cheaper and more numerous the derivatives, the longer the chain and the weaker the average link. This is the single most important unmanaged risk created by the trend this book otherwise commends.

\emph{Agency.} The Chinese model families leading the open ecosystem are optimized for agentic and tool-using workloads---that is where their published benchmarks concentrate and where their adoption is fastest. Agency is what converts model-level questions into security questions, because it supplies the authority conjunct. Government evaluation supports the narrower point that ordinary released models differ widely in their resistance to indirect injection: in one self-hosted evaluation derived from a public agent benchmark, a released open reasoning model attempted credential exfiltration, phishing transmission, and malware download or execution in 37\%, 48\%, and 49\% of applicable trajectories, against corresponding means of 4\%, 3\%, and 4\% for two US frontier models and 27\%, 48\%, and 28\% for a US open-weight model~\cite{caisi,order66}. The metric counts \emph{attempted} malicious actions even where tools prevented completion, and it measures the susceptibility of particular model--agent configurations under tested conditions---not latent backdoors, not developer intent, and not a property of national origin, as the comparable figures for the US open-weight model make plain. What it does establish is that \emph{injection resistance is an engineering property that varies by an order of magnitude between releases, and it is not what open-weight leaderboards measure}.

\emph{Local inference.} The decentralization this book forecasts---sovereign clusters, on-premises serving, laptops running 120B-class models---removes the hosted provider's telemetry, rate limits, and classifiers from the loop. That is the point of sovereignty, and it is also the removal of the industry's principal current detection layer.

\section{From possibility to probability}
\label{sec:riskchain}

The chapter has so far obeyed its evidence hierarchy, but the strategic narrative still risks the move this book criticizes elsewhere: jumping from demonstrated components to systemic consequence in one step. The tempting bridge is a simple product of six probabilities, and it is wrong: a bare product smuggles in an independence assumption exactly where the events are most correlated. The formulation adopted here is a \emph{fault tree}. Systemic loss requires the conjunction of stages, but the first stage is a disjunction of compromise paths:
\begin{align*}
\mathcal{L}_{\mathrm{sys}} \;=\;{}& \bigl(\text{artifact compromise} \lor \text{runtime compromise} \lor \text{supply-chain compromise}\bigr)\\
&\land\; \text{privileged deployment} \;\land\; \text{trigger success}\\
&\land\; \bigl(\text{propagation} \lor \text{correlated activation}\bigr) \;\land\; \text{containment failure}.
\end{align*}
The structure earns its keep in three ways a product cannot. \emph{First, the OR-gates are where the wild incidents actually live}: every real-world event in this chapter's catalogue---the platform breach, the malicious skills, the worm demonstration---entered through the runtime or supply-chain branch, not the weights branch, so a model-only risk analysis prunes the tree at precisely the branches with observed traversals. \emph{Second, correlation enters as common-cause nodes rather than as an afterthought}: monoculture is a single upstream node feeding both ``privileged deployment'' (everyone deploys the same family) and ``propagation/correlated activation'' (a shared flaw fires everywhere at once), and broad agent permissions feed both ``trigger success'' and ``containment failure.'' A fault tree makes the common causes explicit and prices them once, where a product multiplies their effects as if they were separate luck. \emph{Third, every edge now carries a defense ledger}: for each gate, the analysis must list the preventive controls, the detective controls, the recovery controls, the \emph{evidence that each control works}, and the residual uncertainty---which is exactly the content of Appendix~\ref{app:controls}, reorganized by the tree so that a defender can see which conjunct each of the fourteen control families cuts and which gates remain uninstrumented.

Bounding numbers still have their place, but they are too easily generated the lazy way---by multiplying endpoint values, which is an independence calculation wearing a fault tree's clothes. The treatment below states its dependence model explicitly, using the standard common-cause instrument of reliability engineering, a beta-factor decomposition. Write each gate probability as a mixture of an independent component and a common-cause component driven by the shared upstream nodes (monoculture $M$, broad permissions $B$):
\begin{align*}
\Pr(\mathcal{L}_{\mathrm{sys}}) \;\approx\; \underbrace{\textstyle\prod_i p_i}_{\text{independent part}} \;+\; \beta_{\mathcal{M}} \cdot \Pr(\mathcal{M})\cdot \min(p_{\mathrm{deploy}}, p_{\mathrm{prop}}) \;+\; \beta_{\mathcal{B}} \cdot \Pr(\mathcal{B})\cdot \min(p_{\mathrm{trig}}, p_{\mathrm{contain}}),
\end{align*}
Here $p_i$ ranges over the five gate probabilities listed below, in the order of the tree; $p_{\mathrm{deploy}}$, $p_{\mathrm{trig}}$, $p_{\mathrm{prop}}$ and $p_{\mathrm{contain}}$ are the privileged-deployment, trigger-success, propagation and containment-failure gates respectively. $\mathcal{M}$ denotes the event that a monoculture exists deep enough to correlate deployments---so $\Pr(\mathcal{M})$ is not a market share but the probability that concentration reaches the level at which one flaw is simultaneously present across the deployed base, and this book's forecast concentration of 0.5--0.8 is used as a proxy for it, a substitution the reader should treat as the weakest link in the chain. $\mathcal{B}$ denotes the event that agent permissions are broad enough to couple trigger success to containment failure; no independent estimate of $\Pr(\mathcal{B})$ exists, so it is swept over 0.3--0.7 rather than asserted. The coefficients $\beta_{\mathcal{M}}, \beta_{\mathcal{B}} \in [0,1]$ are \emph{beta factors} in the sense of reliability engineering: the fraction of a gate's failure probability attributable to the shared cause rather than to independent chance, with the $\min(\cdot,\cdot)$ form taken because a common cause cannot induce more correlated failure than the smaller of the two gates it couples. With the per-gate ranges stated here [S] (compromise path traversed 0.05--0.5; privileged deployment 0.2--0.6; trigger success 0.05--0.4; propagation or correlated activation 0.05--0.3; containment failure 0.05--0.3): the independent product spans $\sim$$1\times 10^{-6}$ to $\sim$$1\times 10^{-2}$; and the common-cause terms, with $\beta_{\mathcal{M}},\beta_{\mathcal{B}}$ of 0.1--0.3 (the range empirically observed for common-cause failures in redundant engineered systems) and $\Pr(\mathcal{M})$ set to the monoculture concentration this book itself forecasts (0.5--0.8), \emph{add} contributions of order $10^{-3}$ to $10^{-2}$ at the pessimistic end---meaning the correlated tree's upper bound is reached not by every gate failing at its worst value simultaneously, but by two common causes doing most of the work. That is the analytically honest form of the claim: the five-year range remains $\sim$$10^{-6}$ to $\sim$$10^{-2}$, but the upper end is now \emph{driven by the correlation structure}, not by multiplied pessimism, and the single most effective intervention is whichever control cuts a common-cause $\beta_{\mathcal{M}}$ or $\beta_{\mathcal{B}}$---which runtime diversity and deterministic tool mediation do, and model-family diversity alone does not.

That last observation deserves its own table, because it inverts the intuitive diversity policy:

\begin{table}[htbp]
  \centering
  \small
  \caption{Four deployment configurations against the fault tree [S]. ``Common-cause exposure'' names which beta factors remain live; the ranking is the analysis's central policy output.}
  \label{tab:configs}
  \begin{tabular}{p{44mm}p{40mm}p{52mm}}
    \toprule
    Configuration & Common-cause exposure & Systemic-risk reading \\
    \midrule
    Diverse models \emph{and} runtimes, brokered authority & both beta factors cut & lowest; the tree approaches its independent floor \\
    \addlinespace
    Common model, diverse runtimes & $\beta_{\mathcal{M}}$ live at model layer only & moderate: a weight-level flaw is correlated, but propagation paths and containment differ per site \\
    \addlinespace
    Diverse models, common runtime & $\beta_{\mathcal{M}}$ live at runtime layer & \emph{higher than the previous row}: every observed wild incident traversed the runtime/supply-chain branch, and a shared framework correlates exactly those \\
    \addlinespace
    Common model and runtime, broad privileges & both beta factors at maximum & the upper bound of the analysis; approximately the current default deployment pattern \\
    \bottomrule
  \end{tabular}
\end{table}

Table~\ref{tab:configs}'s third row is the uncomfortable one: a nation that diversifies its model portfolio while standardizing on one agent framework has bought \emph{less} risk reduction than one that runs a single model family across isolated, diverse runtimes with deterministic tool mediation---because the runtime branch is where the realized incidents live. The policy content of the whole section sharpens accordingly: \emph{the terms are not equally controllable}. Deployment authority, runtime diversity, and containment are engineering decisions available to any operator today; artifact compromise and trigger discovery are largely not. Defenders should spend on the gates they own and the beta factors they can cut---and the fault tree, unlike the product, shows exactly which those are.

\subsection{Model monoculture is not the only monoculture}

One refinement changes what diversity is worth. Thousands of derivatives of one base model share learned vulnerabilities, but systemic compromise also depends on shared \emph{runtime}: the same inference servers, agent frameworks, package managers, tool schemas, identity providers, and cloud control planes. Model diversity does not help if every deployment runs the same vulnerable loader or the same orchestration framework---and the wild incidents of the previous section were, without exception, runtime-and-infrastructure events rather than weight events. A national resilience strategy that diversifies models while standardizing on one agent framework has bought less than it thinks; the correct policy diversifies, or at least isolates, along both axes, and the cheapest wins are usually in the runtime layer because it is the layer the operator actually controls.

\section{Consequences for the strategy}

Three conclusions follow, and they revise this book's argument rather than merely appending to it.

\textbf{First, the defensible objective is not trustworthy models but unnecessary model trust.} Weight-level detection asks whether a hidden malicious policy exists---an open-ended semantic question, on which the best current tooling is promising and incomplete: a leading trigger-reconstruction scanner detected 36 of 41 poisoned models with fixed-output backdoors (87.8\%) with zero false positives among 13 clean models, but reconstructed a working trigger in most-but-not-all cases for a more variable insecure-code task, on models from 270M to 14B parameters, with no formal guarantee and known blind spots for context-dependent triggers~\cite{haystack,order66}. Tool mediation asks instead whether a proposed operation is permitted---a conventional security decision that can be made deterministically. Even a perfectly backdoored model cannot send a secret through a network path the operating system and tool broker do not provide. The controls that follow are the ones any state adopting the commons must fund: separate inference from action execution; a deterministic tool broker with operation, destination, argument, volume, and rate allow-lists outside the model; separation of secret access from egress; payload inspection rather than text inspection alone; approval gates for high-impact actions; core policy and configuration immutable to the agent; memory writes treated as security decisions; constrained destructive filesystem operations with independent backup credentials; complete external logs in append-only storage unavailable to the model; frozen, signed, hash-pinned tool schemas; a hardened artifact supply chain with controlled mirrors, signed manifests, tensor-only formats, and software bills of materials; evaluation environments contained as hostile-code execution with independently verified egress denial; red-teaming of the actual integration rather than the model in isolation; and no reliance on fine-tuning as sanitization~\cite{order66}.

\textbf{Second, this is the strongest available brake on the adoption curve---and it is legitimate.} Chapter~\ref{ch:analogy} identified certification regimes as the property that has defeated the industrial playbook in aircraft and pharmaceuticals, and Chapter~\ref{ch:synthesis} listed a hardening trust barrier as falsification branch~(f). This chapter supplies its technical content. A regulator who requires signed provenance, reproducible lineage, injection-resistance testing, and brokered authority for agentic deployment in critical sectors is not protecting an incumbent; they are applying to model artifacts the discipline that airframes and drugs already carry. The commons will clear that bar---but clearing it takes institutions, and institutions take years. That lag, not tariffs, is the West's real remaining source of deployment-layer time, and it is available to any jurisdiction willing to build the apparatus rather than merely publish the rule.

\textbf{Third, the certification vacancy is now urgent rather than merely available.} Chapter~\ref{ch:trust} identified the unowned role: whoever operates the evaluation batteries, signs the manifests, and runs the mirrors becomes the certification authority of the open commons. The evidence in this chapter shows what that authority must actually do---maintain reproducible artifact lineage across a derivative chain that no single vendor controls, run injection and backdoor evaluation on the \emph{integration} rather than the model, and publish results that neither bloc's marketing controls. No commercial actor has an incentive to do this honestly for models it does not sell, and no single government will be trusted to do it for models it did not train. That is an argument for an international, institutionally neutral body---which is the subject of Chapter~\ref{ch:consortium}.
